\epigraph{\textit{``Reading maketh a full man; conference a ready man; and writing an exact man.''}}{--- \textit{Francis Bacon}~\citep{bacon1625studies}}
What makes a writer exact has since been studied as a process in which authors coordinate planning, drafting, and revision~\cite{rohman1965pre, flower1981cognitive, hayes1996new, hayes2012modeling}. \citet{flower1981cognitive} made that coordination a matter of control. The processes recur in no fixed order, and a \texttt{Monitor}, guided by goals the writer keeps revising, decides which runs next. That account has guided writing research and the design of tools that support human writers~\cite{gero2022design}, yet no writing system has made process selection its controller. Language agents make such a writer possible, and two questions follow: does a writer improve when its control follows the theory, and which part of the theory accounts for the gain?

\begin{figure*}[t]
\centering
\resizebox{0.95\textwidth}{!}{%
\begin{tikzpicture}[
    font=\fontsize{8.5}{10}\selectfont,
    box/.style={
        draw,
        rounded corners=2pt,
        align=center,
        minimum height=9mm,
        inner sep=4pt
    },
    input/.style={box,text width=27mm},
    state/.style={box,text width=33mm},
    monitor/.style={box,text width=31mm},
    role/.style={box,text width=25mm},
    output/.style={box,text width=27mm},
    arrow/.style={-{Latex[length=1.7mm]},semithick},
    biarrow/.style={<{Latex[length=1.7mm]}-{Latex[length=1.7mm]},semithick}
]

\node[input] (task) at (0,1.6)
    {\textbf{Writing task}\\context + constraints};

\node[state] (workspace) at (3.8,1.6)
    {\textbf{Shared workspace}\\draft\\goal network\\process history};

\node[monitor] (monitor) at (7.6,1.6)
    {\textbf{Host agent: \texttt{Monitor}}\\inspect state\\choose next writing activity};

\node[role] (planner) at (5.1,-0.5)
    {\textbf{\texttt{Planner}}\\planning};
\node[role] (translator) at (7.6,-0.5)
    {\textbf{\texttt{Translator}}\\drafting / translating};
\node[role] (reviewer) at (10.1,-0.5)
    {\textbf{\texttt{Reviewer}}\\reviewing};

\node[output] (output) at (12.1,1.6)
    {\textbf{Completed document}};

\draw[arrow] (task) -- (workspace);
\draw[biarrow] (workspace) -- (monitor);
\draw[arrow] (monitor.south) to[out=-120,in=90] (planner.north);
\draw[arrow] (monitor.south) -- (translator.north);
\draw[arrow] (monitor.south) to[out=-60,in=90] (reviewer.north);

\draw[arrow] (planner.west) -| (workspace.south);
\draw[arrow] (translator.west) to[out=180,in=-55] (workspace.south);
\draw[arrow] (reviewer.west) to[out=180,in=-30] (workspace.south);

\draw[arrow] (monitor) -- node[above]{finish} (output);

\node[align=center,font=\itshape] at (7.6,-1.55)
    {Each delegated activity updates the workspace before control returns to the \texttt{Monitor}.};

\end{tikzpicture}%
}
\caption{Overview of \condAgenticCogWriter. A writing task is carried through a persistent shared workspace. At each iteration, the host-agent \texttt{Monitor} inspects the current draft, goal network, and process history, chooses whether the document next needs planning, drafting, or reviewing, and delegates that activity to the corresponding role agent. The updated state is reconsidered until the \texttt{Monitor} terminates with a completed document.}
\label{fig:agentic-cogwriter}
\end{figure*}


Neither question has been answered: language-model writing systems adapt over units of work such as sections, subtasks, and plans~\cite{bai2025longwriter, xiong2025beyond}, never over the writing processes themselves.

In this paper, we present Agentic CogWriter, a writing agent whose controller chooses the next writing process rather than the next unit of work, and use it to answer both questions. The system is built from skills and subagents on a coding-agent harness: the \texttt{Monitor} is a skill, the \texttt{Planner}, \texttt{Translator}, and \texttt{Reviewer} subagents run the three processes, and the theory's control commitments become ablations that change instructions rather than code. We compare seven writers on the most demanding prompts of the same three benchmarks, with three independent replications. The architecture is favored over the baselines, winning \PooledFourOneRateMean\ of non-tied comparisons against single-pass writing.

The distinction matters because existing systems are not short of structure. They plan before drafting~\cite{yang2022re3, shao2024assisting}, decompose the task or refine a draft in rounds~\cite{bai2025longwriter, wan2025cognitive, madaan2023self}, and interleave planning with execution~\cite{xiong2025beyond}; what none of them does is make the choice of the next writing process the decision its controller takes. \citet{flower1981cognitive} place control exactly there: their \texttt{Monitor} chooses among writing processes, the choice depends on the writer's goals, and evaluating a partial draft can interrupt any process and revise those goals. Cognitively inspired systems borrow parts of this account~\cite{wan2025cognitive}, and the systems above share its vocabulary of planning, drafting, and revision~\cite{yang2022re3, madaan2023self, shao2024assisting, bai2025longwriter, wan2025cognitive}, but none adopts its controller.

Agentic CogWriter delegates \texttt{Planning}, \texttt{Translating}, and \texttt{Reviewing} to separate agents over a shared draft, a goal network, and a history, and the \texttt{Monitor} takes control back after each process. The same plugin runs unchanged on two host platforms, so the writer described here is one a reader can run.

The theory makes three commitments that an implementation can keep or drop~\cite{flower1981cognitive}: the writer maintains a network of goals that it keeps revising, a \texttt{Monitor} chooses which process runs next rather than following a fixed order, and the processes are distinct rather than one undifferentiated act of writing. We turn each commitment into an ablation. The goal-network ablation drops the goal network, the fixed-order ablation replaces the \texttt{Monitor}'s choice with a fixed cycle, and the single-writer ablation keeps one agent writing paragraph by paragraph with no delegation; a loss after removing a commitment would show that it contributed, whereas no detectable loss does not show that it contributed nothing. We compare the full system and these three ablations with three writers built the way current systems are built, a single call, fixed stages, and adaptive task planning, on WritingBench~\cite{wu2026writingbench}, HelloBench~\cite{que2024hellobench}, and DoLoMiTes~\cite{malaviya2025dolomites}, with the comparisons and their correction for multiple testing fixed in advance (Section~\ref{sec:evaluation}); Section~\ref{sec:research-questions} states the three research questions.

This paper makes two contributions. First, it turns the cognitive process theory of writing into a running writer: the \texttt{Monitor} is a skill on a coding agent, the \texttt{Planner}, \texttt{Translator}, and \texttt{Reviewer} subagents run the three processes, and each of the theory's commitments can be switched off by changing an instruction. Second, it reports which of those commitments made a measurable difference. Only one did: when the three processes are folded back into one agent, the writer loses clearly, whereas removing the goal network or the dynamic process order shows no detectable loss, and the traces show the \texttt{Monitor} mostly following the textbook order on its own. Two caveats apply throughout. The single-agent ablation also changes how finely the writer recurses and how much it writes, so the loss cannot be attributed to delegation alone; and every confirmatory judgment came from a model in the generator's family, which we check with a judge from another family and with output-length matching.
